\documentclass[a4paper,12pt]{article}
\usepackage[italian]{babel}
\usepackage[utf8]{inputenc}
\usepackage[T1]{fontenc}
\usepackage{geometry}
\usepackage{setspace}
\usepackage{enumitem}
\usepackage{listings}
\usepackage{xcolor}
\geometry{margin=2.5cm}
\onehalfspacing

\definecolor{codegray}{rgb}{0.95,0.95,0.95}
\lstset{
    backgroundcolor=\color{codegray},
    basicstyle=\ttfamily\small,
    breaklines=true,
    frame=single,
    showstringspaces=false,
    tabsize=4,
    language=SQL
}

\title{25 query SQL di difficoltà crescente - Soluzioni}
\author{Prof. Fedeli Massimo}
\date{}

\begin{document}

\maketitle

\section*{Schema logico di riferimento}

\begin{flushleft}
STUDENTE(\underline{Id\_studente}, Cognome, Nome, Data\_nascita, Luogo\_nascita, CF, citta, CAP, via, n\_civico, Email\_istituzionale, Classe, Indirizzo\_studi) \\
AZIENDA(\underline{P\_iva}, Ragione\_sociale, Sede\_legale, Sede\_operativa, settore\_attivita, Email\_aziendale, Telefono\_ref) \\
TUTOR\_SCUOLA(\underline{Codice\_docente}, Cognome, Nome, Email\_istituzionale, Dipartimento\_app, N\_max\_studenti\_seguiti) \\
TUTOR\_AZIENDA(\underline{Codice\_dipendente}, Cognome, Nome, Email\_aziendale, P\_iva) \\
COMPETENZE(\underline{nome\_competenza}) \\
DISPONIBILITA(\underline{Codice\_id}, Data\_inizio, Data\_fine, N\_max\_studenti\_ospitabili, Descrizione\_attivita, Indirizzo\_studi\_consigliato) \\
AZIENDA\_DISPONIBILITA(\underline{P\_iva, Codice\_id}) \\
RICHIEDE(\underline{nome\_competenza, Codice\_id}) \\
ESPERIENZA(\underline{Id\_studente, P\_iva\_azienda\_ospitante, Codice\_docente, Codice\_dipendente, Codice\_id\_disponibilita}, Periodo\_effettivo, N\_ore\_previste, N\_ore\_svolte, Valutazione\_finale, esito)
\end{flushleft}


\section*{Soluzioni}\subsection*{Query 1}Elencare cognome, nome, classe ed email istituzionale di tutti gli studenti, ordinati per cognome e nome.
\begin{lstlisting}
SELECT Cognome, Nome, Classe, Email_istituzionale
FROM STUDENTE
ORDER BY Cognome, Nome;
\end{lstlisting}
\subsection*{Query 2}Visualizzare ragione sociale, settore di attività e sede operativa di tutte le aziende con sede operativa a Pesaro.
\begin{lstlisting}
SELECT Ragione_sociale, settore_attivita, Sede_operativa
FROM AZIENDA
WHERE Sede_operativa = 'Pesaro';
\end{lstlisting}
\subsection*{Query 3}Mostrare i tutor scolastici che possono seguire almeno 11 studenti.
\begin{lstlisting}
SELECT Codice_docente, Cognome, Nome, N_max_studenti_seguiti
FROM TUTOR_SCUOLA
WHERE N_max_studenti_seguiti >= 11;
\end{lstlisting}
\subsection*{Query 4}Elencare le disponibilità con Data\_fine successiva al 31 maggio 2026, mostrando codice, descrizione attività e numero massimo di studenti ospitabili.
\begin{lstlisting}
SELECT Codice_id, Descrizione_attivita, N_max_studenti_ospitabili, Data_fine
FROM DISPONIBILITA
WHERE Data_fine > '2026-05-31';
\end{lstlisting}
\subsection*{Query 5}Mostrare gli studenti della classe 5A nati dopo il 1 gennaio 2005.
\begin{lstlisting}
SELECT Id_studente, Cognome, Nome, Data_nascita
FROM STUDENTE
WHERE Classe = '5A' AND Data_nascita > '2005-01-01';
\end{lstlisting}
\subsection*{Query 6}Elencare per ogni esperienza lo studente, l'azienda ospitante e le ore previste.
\begin{lstlisting}
SELECT s.Cognome, s.Nome, a.Ragione_sociale, e.N_ore_previste
FROM ESPERIENZA e
JOIN STUDENTE s ON e.Id_studente = s.Id_studente
JOIN AZIENDA a ON e.P_iva_azienda_ospitante = a.P_iva
ORDER BY s.Cognome, s.Nome;
\end{lstlisting}
\subsection*{Query 7}Mostrare gli studenti che hanno svolto meno di 100 ore effettive, indicando anche azienda ospitante e ore svolte.
\begin{lstlisting}
SELECT s.Cognome, s.Nome, a.Ragione_sociale, e.N_ore_svolte
FROM ESPERIENZA e
JOIN STUDENTE s ON e.Id_studente = s.Id_studente
JOIN AZIENDA a ON e.P_iva_azienda_ospitante = a.P_iva
WHERE e.N_ore_svolte < 100;
\end{lstlisting}
\subsection*{Query 8}Elencare le aziende con la relativa disponibilità pubblicata, mostrando ragione sociale, descrizione attività e indirizzo di studi consigliato.
\begin{lstlisting}
SELECT a.Ragione_sociale, d.Descrizione_attivita, d.Indirizzo_studi_consigliato
FROM AZIENDA a
JOIN AZIENDA_DISPONIBILITA ad ON a.P_iva = ad.P_iva
JOIN DISPONIBILITA d ON ad.Codice_id = d.Codice_id
ORDER BY a.Ragione_sociale;
\end{lstlisting}
\subsection*{Query 9}Visualizzare per ogni disponibilità le competenze richieste, mostrando codice disponibilità, descrizione attività e nome competenza.
\begin{lstlisting}
SELECT d.Codice_id, d.Descrizione_attivita, r.nome_competenza
FROM DISPONIBILITA d
JOIN RICHIEDE r ON d.Codice_id = r.Codice_id
ORDER BY d.Codice_id, r.nome_competenza;
\end{lstlisting}
\subsection*{Query 10}Mostrare per ogni esperienza lo studente, il tutor scolastico e il tutor aziendale associati.
\begin{lstlisting}
SELECT s.Cognome AS Cognome_studente,
       s.Nome AS Nome_studente,
       ts.Cognome AS Cognome_tutor_scuola,
       ts.Nome AS Nome_tutor_scuola,
       ta.Cognome AS Cognome_tutor_azienda,
       ta.Nome AS Nome_tutor_azienda
FROM ESPERIENZA e
JOIN STUDENTE s ON e.Id_studente = s.Id_studente
JOIN TUTOR_SCUOLA ts ON e.Codice_docente = ts.Codice_docente
JOIN TUTOR_AZIENDA ta ON e.Codice_dipendente = ta.Codice_dipendente
ORDER BY s.Cognome, s.Nome;
\end{lstlisting}
\subsection*{Query 11}Calcolare quante aziende appartengono a ciascun settore di attività.
\begin{lstlisting}
SELECT settore_attivita, COUNT(*) AS numero_aziende
FROM AZIENDA
GROUP BY settore_attivita
ORDER BY numero_aziende DESC, settore_attivita;
\end{lstlisting}
\subsection*{Query 12}Calcolare il numero di studenti per ciascuna classe.
\begin{lstlisting}
SELECT Classe, COUNT(*) AS numero_studenti
FROM STUDENTE
GROUP BY Classe
ORDER BY Classe;
\end{lstlisting}
\subsection*{Query 13}Per ogni azienda ospitante, calcolare il numero di esperienze svolte.
\begin{lstlisting}
SELECT a.Ragione_sociale, COUNT(*) AS numero_esperienze
FROM AZIENDA a
JOIN ESPERIENZA e ON a.P_iva = e.P_iva_azienda_ospitante
GROUP BY a.P_iva, a.Ragione_sociale
ORDER BY numero_esperienze DESC, a.Ragione_sociale;
\end{lstlisting}
\subsection*{Query 14}Determinare, per ogni esperienza, il deficit di ore rispetto a quelle previste, mostrando studente, ore previste, ore svolte e differenza.
\begin{lstlisting}
SELECT s.Cognome, s.Nome,
       e.N_ore_previste,
       e.N_ore_svolte,
       (e.N_ore_previste - e.N_ore_svolte) AS ore_mancanti
FROM ESPERIENZA e
JOIN STUDENTE s ON e.Id_studente = s.Id_studente
ORDER BY ore_mancanti DESC, s.Cognome, s.Nome;
\end{lstlisting}
\subsection*{Query 15}Calcolare la media delle ore svolte per ciascun esito finale.
\begin{lstlisting}
SELECT esito, AVG(N_ore_svolte) AS media_ore_svolte
FROM ESPERIENZA
GROUP BY esito;
\end{lstlisting}
\subsection*{Query 16}Trovare gli studenti che hanno svolto un'esperienza presso un'azienda del settore 'Software'.
\begin{lstlisting}
SELECT s.Cognome, s.Nome, a.Ragione_sociale
FROM STUDENTE s
JOIN ESPERIENZA e ON s.Id_studente = e.Id_studente
JOIN AZIENDA a ON e.P_iva_azienda_ospitante = a.P_iva
WHERE a.settore_attivita = 'Software'
ORDER BY s.Cognome, s.Nome;
\end{lstlisting}
\subsection*{Query 17}Elencare le aziende che non appartengono ai settori 'IT' e 'Software'.
\begin{lstlisting}
SELECT Ragione_sociale, settore_attivita
FROM AZIENDA
WHERE settore_attivita NOT IN ('IT', 'Software')
ORDER BY Ragione_sociale;
\end{lstlisting}
\subsection*{Query 18}Mostrare le disponibilità il cui numero massimo di studenti ospitabili è maggiore della media complessiva.
\begin{lstlisting}
SELECT Codice_id, Descrizione_attivita, N_max_studenti_ospitabili
FROM DISPONIBILITA
WHERE N_max_studenti_ospitabili > (
    SELECT AVG(N_max_studenti_ospitabili)
    FROM DISPONIBILITA
)
ORDER BY N_max_studenti_ospitabili DESC, Codice_id;
\end{lstlisting}
\subsection*{Query 19}Trovare gli studenti che hanno svolto un numero di ore superiore alla media delle ore svolte da tutti gli studenti.
\begin{lstlisting}
SELECT s.Cognome, s.Nome, e.N_ore_svolte
FROM STUDENTE s
JOIN ESPERIENZA e ON s.Id_studente = e.Id_studente
WHERE e.N_ore_svolte > (
    SELECT AVG(N_ore_svolte)
    FROM ESPERIENZA
)
ORDER BY e.N_ore_svolte DESC, s.Cognome, s.Nome;
\end{lstlisting}
\subsection*{Query 20}Elencare i tutor scolastici che non risultano assegnati ad alcuna esperienza.
\begin{lstlisting}
SELECT ts.Codice_docente, ts.Cognome, ts.Nome
FROM TUTOR_SCUOLA ts
LEFT JOIN ESPERIENZA e ON ts.Codice_docente = e.Codice_docente
WHERE e.Codice_docente IS NULL
ORDER BY ts.Cognome, ts.Nome;
\end{lstlisting}
\subsection*{Query 21}Determinare il settore di attività con il maggior numero di aziende.
\begin{lstlisting}
SELECT settore_attivita, COUNT(*) AS numero_aziende
FROM AZIENDA
GROUP BY settore_attivita
HAVING COUNT(*) = (
    SELECT MAX(t.numero_aziende)
    FROM (
        SELECT COUNT(*) AS numero_aziende
        FROM AZIENDA
        GROUP BY settore_attivita
    ) t
);
\end{lstlisting}
\subsection*{Query 22}Trovare l'azienda o le aziende che ospitano studenti con il numero massimo di ore svolte.
\begin{lstlisting}
SELECT a.Ragione_sociale, s.Cognome, s.Nome, e.N_ore_svolte
FROM ESPERIENZA e
JOIN AZIENDA a ON e.P_iva_azienda_ospitante = a.P_iva
JOIN STUDENTE s ON e.Id_studente = s.Id_studente
WHERE e.N_ore_svolte = (
    SELECT MAX(N_ore_svolte)
    FROM ESPERIENZA
)
ORDER BY a.Ragione_sociale, s.Cognome, s.Nome;
\end{lstlisting}
\subsection*{Query 23}Per ogni settore di attività, calcolare la media delle ore svolte dagli studenti ospitati dalle aziende di quel settore.
\begin{lstlisting}
SELECT a.settore_attivita, AVG(e.N_ore_svolte) AS media_ore_svolte
FROM AZIENDA a
JOIN ESPERIENZA e ON a.P_iva = e.P_iva_azienda_ospitante
GROUP BY a.settore_attivita
ORDER BY media_ore_svolte DESC;
\end{lstlisting}
\subsection*{Query 24}Elencare le disponibilità associate ad aziende del settore 'AI' oppure 'Security', mostrando azienda, attività e competenza richiesta.
\begin{lstlisting}
SELECT a.Ragione_sociale, d.Descrizione_attivita, r.nome_competenza
FROM AZIENDA a
JOIN AZIENDA_DISPONIBILITA ad ON a.P_iva = ad.P_iva
JOIN DISPONIBILITA d ON ad.Codice_id = d.Codice_id
JOIN RICHIEDE r ON d.Codice_id = r.Codice_id
WHERE a.settore_attivita IN ('AI', 'Security')
ORDER BY a.Ragione_sociale, d.Descrizione_attivita;
\end{lstlisting}
\subsection*{Query 25}Trovare gli studenti che hanno svolto meno ore del corrispondente valore medio delle esperienze della loro stessa classe.
\begin{lstlisting}
SELECT s.Cognome, s.Nome, s.Classe, e.N_ore_svolte
FROM STUDENTE s
JOIN ESPERIENZA e ON s.Id_studente = e.Id_studente
WHERE e.N_ore_svolte < (
    SELECT AVG(e2.N_ore_svolte)
    FROM ESPERIENZA e2
    JOIN STUDENTE s2 ON e2.Id_studente = s2.Id_studente
    WHERE s2.Classe = s.Classe
)
ORDER BY s.Classe, e.N_ore_svolte, s.Cognome, s.Nome;
\end{lstlisting}

\end{document}
